\section{The Language \OurLang}
\label{sec:language}

% PAR: section overview, call-by-value lambda calculus with session channels | example: - | figure/rule: fig:syntax | map: syn-terms

\subsection{Syntax and Semantics}
\label{sec:syntax}

% PAR: ends instead of channels, c+ and c-, dual end, the scheduler names both ends of its six channels | example: ex-milner | figure/rule: fig:syntax | map: syn-ends
% Syntax of GV-prec. From map blocks syn-ends, syn-prims, syn-terms,
% syn-values, syn-let, syn-procs. Named variables in place of de Bruijn
% indices (PLAN.md, presentation).
\begin{figure}[tb]
  \small
  \judghead{Ends, constants, terms and values}{\tm}
  \[
  \begin{array}{@{}r@{\ \ }c@{\ \ }l@{\qquad}r@{}}
    \eo{} & ::= & \ch{}^{+} \mid \ch{}^{-} & \text{ends, } \dual{\ch{}^{\pm}} = \ch{}^{\mp}\\[2pt]
    \cst{} & ::= & \tsend \mid \trecv \mid \tclose \mid \tselect_i \mid \tbranch \mid \tlink \mid \tspawn \mid \tforce & \text{constants}\\[2pt]
    \tm & ::= & \vr{} \mid \tunit \mid \eo{} \mid \tlam{\vr{}}{\tm} \mid \tm_1\,\tm_2 \mid \tpair{\tm_1}{\tm_2}
              \mid \tletu{\tm_1}{\tm_2} \mid \tletp{\vr{}}{\vrB{}}{\tm_1}{\tm_2} & \text{terms}\\
      & \mid & \tinj_i\,\tm \mid \tcasee{\tm}{\vr{}}{\tm_1}{\vr{}}{\tm_2} \mid \tbox{\tmB{}}{\tm} \mid \tnew \mid \cst{}\,\tm & \\[2pt]
    \val & ::= & \tunit \mid \eo{} \mid \tlam{\vr{}}{\tm} \mid \tbox{\tmB{}}{\tm} \mid \tpair{\val_1}{\val_2} \mid \tinj_i\,\val & \text{values}\\[2pt]
    \proc & ::= & (C, E) & \text{processes}
  \end{array}
  \]
  \[ \tletv{\vr{}}{\tm_1}{\tm_2} \;\triangleq\; \tletp{\vr{}}{\valB{}}{\tpair{\tm_1}{\tunit}}{\tletu{\valB{}}{\tm_2}} \qquad (\valB{} \text{ fresh}) \]
  \figrule
  \Description{Grammar of ends, constants, terms, values and processes.}
  \caption{Syntax of \OurLang. A process $\proc = (C, E)$ has a set $C$ of restricted channels and a finite map $E$ from thread names $t$ to terms, the threads $\proc(t) = E(t)$. Term steps $e \xrightarrow{\ell} e'$ and process steps $\proc \xrightarrow{ev} \proc'$ are standard (\cref{app:sem-labels,app:sem-head,app:sem-ctx,app:sem-proc,app:syn-csub}).}
  \label{fig:syntax}
\end{figure}


% PAR: constants k e, thirteen term formers, new as a term, select and branch in the core | example: - | figure/rule: fig:syntax | map: syn-prims, syn-terms
% PAR: labelled semantics standard GV, pointer to appendix | example: - | figure/rule: app:semantics | map: sem-labels, sem-head, sem-ctx, sem-proc
% PAR: processes (C, E), restricted channels, threads with names, the scheduler after its creations and spawns as a process with six restricted channels and seven threads | example: ex-milner | figure/rule: fig:syntax, def:runs | map: syn-procs, sem-events
% PAR: smallest process, two threads close one channel, listing in the appendix | example: ex-close | figure/rule: app:ex-close | map: syn-procs, sem-proc
% PAR: events name the threads of a step, closed steps, runs, maximal and weakly fair runs, pending threads | example: - | figure/rule: def:runs | map: sem-events
\begin{definition}[Steps and runs \rocq{istep} \rocq{maximal} \rocq{wfair}]
\label{def:runs}
Processes step as in \cref{app:sem-proc}.
\begin{itemize}
  \item A step $\proc \xrightarrow{ev} \proc'$ has an event $ev$, which records the threads $\mathsf{parts}(ev)$ that take part and their labels (\cref{app:sem-events}).
  \item The event is internal unless a thread offers an action on a free channel to the environment. A closed step $\proc \to \proc'$ has an internal event.
  \item A run of $\proc$ is a sequence $\proc_0 \proc_1 \ldots$ with $\proc_0 = \proc$ and, at each $i$, either a closed step $\proc_i \xrightarrow{ev_i} \proc_{i+1}$ or $\proc_{i+1} = \proc_i$ with no event.
  \item A run is maximal when $\proc_{i+1} = \proc_i$ only where $\proc_i$ has no closed step.
  \item A run is weakly fair when each thread $t$ at each $i$ has some $j \geq i$ at which $t$ takes part in $ev_j$ or takes part in no closed step of $\proc_j$.
  \item A thread $t$ is pending in $\proc$ when the next step of $\proc(t)$ is a session action, a link included.
\end{itemize}
$(C, \{\tm_1, \ldots, \tm_n\})$ is the process whose threads $\tm_1, \ldots, \tm_n$ have distinct names. Typing and the relations of \cref{sec:typing,sec:relation} read only the restricted channels and the terms of a process, not the names of its threads.
\end{definition}

% PAR: terminated, deadlock free, the scheduler with one agent reversed waits in a ring | example: ex-milner | figure/rule: def:df | map: sem-props
% PAR: lock freedom, every pending thread eventually takes part, weak fairness, deadlock freedom lets a thread wait forever while others step | example: - | figure/rule: def:df | map: sem-props
\begin{definition}[Deadlock freedom and lock freedom \rocq{deadlock_free} \rocq{lock_free}]
\label{def:df}\label{def:lf}
\begin{align*}
  \terminated\,\proc &\iff \exists \val.\ \proc = (\varnothing, \{\val, \tunit, \ldots, \tunit\}) \\
  \dfree\,\proc &\iff \forall \proc'.\ \proc \to^{*} \proc' \Rightarrow \terminated\,\proc' \lor \exists \procB.\ \proc' \to \procB \\
  \lfree\,\proc &\iff \text{each maximal weakly fair run } \proc_0 \proc_1 \ldots \text{ of } \proc \text{ has} \\
  &\qquad \forall i, t.\ t \text{ pending in } \proc_i \Rightarrow \exists j \geq i.\ t \in \mathsf{parts}(ev_j)
\end{align*}
Runs, pending threads and $\mathsf{parts}$ are those of \cref{def:runs} (\cref{app:sem-events,app:sem-props}).
\end{definition}
% GAP: sem-props shows the deadlocking reversed cycle only in its breaks text, no map example with a listing for a deadlocked process (the scheduler with A_1 reversed is a candidate)

\subsection{Points and Worlds}
\label{sec:points}

% PAR: point of an end, where it acts next, terms never mention points, the points of the scheduler's ends | example: ex-milner | figure/rule: fig:points-example | map: syn-points
% (a) the points and facts of map example ex-cycle (its code comments).
% (b) where a payload end arrives, from syn-moves and E-Send/E-Recv, with
% the shape of ex-ring (a created at the field x.0 and sent on y at x).
\begin{figure}[tb]
  \centering
  \begin{subfigure}[b]{0.5\linewidth}
    \centering
    \begin{tikzpicture}[x=1cm, y=1cm]
      \node[enddot, label={[endlabel]left:$c : \rho$}] (c) at (0, 0.7) {};
      \node[enddot, label={[endlabel]left:$d : \fld{\rho}{0}$}] (d) at (0, -0.7) {};
      \node[enddot, label={[endlabel]right:$c' : \dual{\rho}$}] (c') at (3.4, 0.7) {};
      \node[enddot, label={[endlabel]right:$d' : \dual{\fld{\rho}{0}}$}] (d') at (3.4, -0.7) {};
      \begin{scope}[on background layer]
        \node[thread, fit=(c)(d), minimum width=1.9cm, xshift=-0.55cm, label={[threadlabel]above:main}] {};
        \node[thread, fit=(c')(d'), minimum width=1.9cm, xshift=0.55cm, label={[threadlabel]above:spawned}] {};
      \end{scope}
      \draw[chan] (c) -- (c');
      \draw[chan] (d) -- (d');
      \draw[fact] (c) -- node[factlabel, pos=0.22, above, sloped] {$\rho \prec \dual{\fld{\rho}{0}}$} (d');
      \draw[fact] (c') -- node[factlabel, pos=0.22, above, sloped] {$\dual{\rho} \prec \fld{\rho}{0}$} (d);
    \end{tikzpicture}
    \caption{Cycle (\cref{ex:cycle})}
    \label{fig:points-example:cycle}
  \end{subfigure}\hfill
  \begin{subfigure}[b]{0.47\linewidth}
    \centering
    \begin{tikzpicture}[x=1cm, y=1cm]
      \node[enddot, label={[endlabel]left:$y : x$}] (y) at (0, 0.7) {};
      \node[enddot, label={[endlabel]left:$a : \fld{x}{0}$}] (a) at (0, -0.7) {};
      \node[enddot, label={[endlabel]right:$\dual{y} : \dual{x}$}] (yb) at (2.7, 0.7) {};
      \node[ghostdot, label={[endlabel]right:$a : \fld{x}{0}$}] (a2) at (2.7, -0.7) {};
      \begin{scope}[on background layer]
        \node[thread, fit=(y)(a), minimum width=1.9cm, xshift=-0.55cm, label={[threadlabel]above:sender}] {};
        \node[thread, fit=(yb)(a2), minimum width=1.9cm, xshift=0.55cm, label={[threadlabel]above:receiver}] {};
      \end{scope}
      \draw[chan] (y) -- (yb);
      \draw[msg] (a) to[bend right=18] node[factlabel, below] {$\tsend\,\tpair{a}{y}$ at $x$} (a2);
    \end{tikzpicture}
    \caption{A payload end arrives at $\fld{x}{0} = \fld{\dual{x}}{0}$}
    \label{fig:points-example:payload}
  \end{subfigure}
  \caption{Points of ends and facts of a world}
  \label{fig:points-example}
\end{figure}


% PAR: an action moves an end from x to x+, choice to a field, dual point | example: ex-milner | figure/rule: fig:worlds | map: syn-moves
% PAR: payload points, the j-th payload end arrives at x.j, both ends compute the same field | example: ex-ring | figure/rule: fig:points-example | map: syn-moves, syn-points, ty-payload
% GAP: no map example in the deadlock-free fragment sends an end and states all its facts. Panel (b) of fig:points-example is schematic, from syn-moves and the shape of ex-ring.
% Points, moves, worlds. From map blocks syn-labels, syn-points, syn-moves,
% w-world, w-order, w-wf, w-ext. Restriction and the join lemma (w-ext) are
% in the appendix. Point constants are written rho (map: a).
\begin{figure}[tb]
  \small
  \judghead{Points}{x}
  \[
  \begin{array}{@{}c@{}}
    b ::= \rho \mid \iota \qquad
    x, y, z ::= b \mid \pplus{x} \mid \fld{x}{\kappa} \mid \dual{x} \qquad (\kappa \in \mathbb{N})
    \qquad\quad
    \dual{\dual{x}} = x \qquad \fld{\dual{x}}{\kappa} = \fld{x}{\kappa} \qquad \dual{\pplus{x}} = \pplus{\dual{x}}
  \end{array}
  \]
  \judghead{Where an end moves, where a payload arrives}{x \mapsto x'}
  \[
  \begin{array}{@{}l@{\qquad}l@{\qquad}l@{\qquad}l@{}}
    x \mapsto \pplus{x} & x \mapsto \dual{\fld{x}{i}} & x \mapsto \fld{x}{i} & \fld{x}{j} = \fld{\dual{x}}{j} \\
    \text{\footnotesize send, receive, close at $x$} & \text{\footnotesize $\tselect_i$ at $x$} & \text{\footnotesize $\tbranch$ at $x$, side $i$} & \text{\footnotesize $j$-th payload end of a message at $x$}
  \end{array}
  \]
  \judghead{Worlds}{W}
  \[
  \begin{array}{@{}l@{\qquad\quad}l@{}}
    W = (A, {\prec_A}),\ \ {\prec_A} \subseteq \Point_A^2 &
    {\ordW} = \mathsf{tc}\big({\prec_A} \cup \{x \prec \pplus{x}\} \cup \{x \prec \pplus{\dual{x}}\}\big) \\[3pt]
    \cons\,W \iff \forall x.\ x \not\ordW x &
    x \preccurlyeq_W y \iff x = y \lor x \ordW y \\[3pt]
    W \wext W' \iff A \subseteq A' \land {\ordW} \subseteq {\ordW[W']} &
    W_1 \wjoin W_2 = (A_1 \cup A_2,\ {\prec_{A_1}} \cup {\prec_{A_2}}) \\[3pt]
    W \grow{\rho}{C} = (A \cup \{\rho\},\ {\prec_A} \cup C[\rho]) &
  \end{array}
  \]
  \figrule
  \Description{Grammar of points, the moves of an end, and the definitions of worlds, their order, consistency, extension, growth and join.}
  \caption{Points, moves and worlds. $\Point_A$ is the set of points with a base in $A$, $\mathsf{tc}$ the transitive closure (restriction and the join lemma in \cref{app:w-ext}).}
  \label{fig:worlds}
\end{figure}


% PAR: world, bases and facts, order generated with program order | example: ex-milner | figure/rule: fig:worlds | map: w-world, w-order
% PAR: the scheduler's world, nine facts and their shifts, all rise with the field label | example: ex-milner | figure/rule: fig:world-example | map: w-world, r-new
% The world of map example ex-milner: the nine facts of ms_base in
% theories/examples/MilnerSched.v, without their shifts x+ < y (ms_L).
% Columns are field labels, the top row holds the points rho.k, the
% bottom row their duals. Arrow styles tell the agent that needs the fact.
\begin{figure}[tb]
  \centering
  \begin{tikzpicture}[x=1.95cm, y=1cm]
    \foreach \k/\pt/\tn/\bn in {%
        0/{\rho}/c_0/d_1, 1/{\fld{\rho}{1}}/a_1/b_1, 2/{\fld{\rho}{2}}/c_1/d_2,%
        3/{\fld{\rho}{3}}/a_2/b_2, 4/{\fld{\rho}{4}}/c_2/d_0, 5/{\fld{\rho}{5}}/a_0/b_0} {
      \node[ptnode] (t\k) at (\k, 1.3) {$\pt$};
      \node[ptnode] (b\k) at (\k, 0) {$\dual{\pt}$};
      \node[endname, above=0pt of t\k] {$\tn$};
      \node[endname, below=0pt of b\k] {$\bn$};
    }
    % A_1
    \draw[factA] (b0) -- (b1);
    \draw[factA] (b0) to[bend right=38] (b2);
    \draw[factA] (t1) -- (b2);
    % A_2
    \draw[factB] (b2) -- (b3);
    \draw[factB] (b2) to[bend right=38] (b4);
    \draw[factB] (t3) -- (b4);
    % A_0
    \draw[factC] (t0.north east) .. controls (1.0, 2.45) and (3.4, 2.45) .. (t4.north west);
    \draw[factC] (t0.north east) .. controls (0.8, 3.0) and (6.7, 3.0) .. (b5.east);
    \draw[factC] (b4) -- (b5);
    % legend
    \begin{scope}[shift={(0, -1.75)}]
      \draw[factA] (0, 0) -- (0.3, 0); \node[legend] at (0.35, 0) {needed by $A_1$};
      \draw[factB] (1.5, 0) -- (1.8, 0); \node[legend] at (1.85, 0) {needed by $A_2$};
      \draw[factC] (3, 0) -- (3.3, 0); \node[legend] at (3.35, 0) {needed by $A_0$};
    \end{scope}
  \end{tikzpicture}
  \Description{The twelve points of the scheduler in six columns, one per field label, with the points in the top row and their duals in the bottom row. Nine arrows for the facts all point from a lower field label to a higher one.}
  \caption{The world of the scheduler (\cref{ex:milner}). An arrow from $x$ to $y$ is a fact $x \prec y$ of the first creation, the grey name next to a point is the end at that point, and every fact also holds from $\pplus{x}$ (not drawn).}
  \label{fig:world-example}
\end{figure}


% PAR: consistent world, no cycle, the scheduler's world consistent because the facts rise | example: ex-milner | figure/rule: fig:world-example | map: w-wf
% PAR: extension, growth by a creation, join, restriction | example: ex-milner | figure/rule: fig:worlds | map: w-ext
% PAR: mechanization insight, one notion of world in the paper, two in the Rocq (static worlds of typing and worlds of the relation) | example: - | figure/rule: fig:worlds | map: w-world, w-ext

\FloatBarrier
